\documentclass[10pt]{article}
\usepackage[T1]{fontenc}
\usepackage{lmodern}
\usepackage[margin=0.8in]{geometry}
\usepackage{amsmath,amssymb,amsthm}
\usepackage[hidelinks]{hyperref}
\setlength{\parindent}{0pt}
\setlength{\parskip}{4pt}
\newtheorem{theorem}{Theorem}
\newcommand{\C}{\mathbb C}
\newcommand{\N}{\mathbb N}
\newcommand{\A}{\mathcal A}
\begin{document}
\begin{center}
{\Large A third-moment obstruction to the claimed free CLT rate}\\[4pt]
{\large Disproof of the universal moment-rate clause in TLMC 00000007730}\\[3pt]
\texttt{lizaixi01}\qquad 24 October 2026
\end{center}

\textbf{Exact scope.}
The source claims an $O(n^{-1})$ moment-convergence rate for centered
self-adjoint tuples free over $B$ with bounded fourth moments. Neither
language assumes symmetry or a vanishing third moment. We use the
scalar specialization $B=\C$ and one-entry tuples. Only this universal
rate clause is refuted. The complete bilingual statement is preserved
in \texttt{SOURCE.md}. We make no assertion about a theorem restricted
to even moments, or the later fourth-cumulant and colored-Motzkin clauses.

\begin{theorem}
There is a concrete unital star-algebra $\A$ of bounded operators on a
complex Hilbert space, with a positive unital bounded tracial state $\tau$,
and a freely independent identically distributed sequence of bounded
self-adjoint elements $Y_i\in\A$ such that
\[
 \tau(Y_i)=0,\quad \tau(Y_i^2)=1,\quad \tau(Y_i^4)=\frac32,
 \qquad
 S_n=\frac1{\sqrt n}\sum_{i=0}^{n-1}Y_i.
\]
For every positive integer $n$,
\[
 \tau(S_n^3)=\frac1{\sqrt{2n}},\qquad
 \int x^3\,d\mu_{\rm sc}(x)=0,
\]
where $\mu_{\rm sc}$ is the actual standard semicircle probability measure.
The discrepancy is not $O(n^{-1})$ as $n\to\infty$.
\end{theorem}

\textbf{1. The actual free product, Hilbert space and trace.}
Let $F_i=C_3=\langle g_i\mid g_i^3=1\rangle$ and
$G=*_{i\in\N}F_i$. Lean uses the native countable indexed group coproduct
and its genuine factor embeddings. Its reduced-word equivalence proves
the normal-form theorem; countability follows from the actual list of
factor letters. Set $H=\ell^2(G;\C)$, using native \texttt{lp} with exponent
$2$. Actual coordinate reindexing gives linear isometry equivalences
\[
 (\lambda(g)f)(h)=f(g^{-1}h).
\]
Their bounded operators satisfy $\lambda(g)\lambda(h)=\lambda(gh)$
and $\lambda(g)^*=\lambda(g^{-1})$, with both unitary identities proved.

The universal group-algebra map is the actual algebra homomorphism
$\rho:\C[G]\to\mathcal B(H)$ obtained from these operators. For
$\delta=\delta_e$, Lean proves the coefficient identity
\[
 (\rho(a)\delta)(h)=a_h.
\]
Thus $\rho$ is faithful. Its image is closed under the genuine adjoint
and forms a native star-subalgebra $\A$ of $\mathcal B(H)$.
This is the concrete unital star-algebra of bounded operators used here,
in the standard algebraic formulation of noncommutative probability
\cite[p.~332]{Novak}.

On this actual algebra define $\tau(T)=\langle\delta,T\delta\rangle$.
The formal proof establishes
\[
 \tau(1)=1,\quad |\tau(T)|\le\|T\|,\quad
 \tau(T^*T)=\|T\delta\|^2,\quad \tau(AB)=\tau(BA).
\]
For traciality, $\tau(\rho(a))=a_e$ and the equality
$(ab)_e=(ba)_e$ is proved for all finitely supported group-algebra
elements by their true linear induction principle. All tracial claims
concern $\A$. The state is bundled as a genuine continuous linear map.

\textbf{2. Full freeness.}
Let $\A_i$ be the actual image of $\C[C_3]$ under the $i$th factor
representation, regarded as a native unital star-subalgebra of $\A$.
If $a\in\C[C_3]$ is centered, its identity coefficient is zero.
Every element in the support of its embedded image therefore comes
from a nonidentity letter of that factor. For an arbitrary product of
such elements, the native support-of-product inclusion extracts an
actual list of nonidentity factor letters. When consecutive factor
indices differ, this list is a nonempty reduced word, so its product
cannot equal $e$ by the native normal-form equivalence. Consequently
\[
 \tau(a_1\cdots a_k)=0
 \quad(k>0,\ a_j\in\A_{i_j},\ \tau(a_j)=0,\ i_j\ne i_{j+1}).
\]
This is proved for every length and every centered element of each
whole factor algebra. It is then transported to the actual represented
star-algebra. The group-algebra coefficient-state construction and this
standard freeness criterion are described in \cite[pp.~339--340]{Novak};
the particular $C_3$ calculations below are proved directly.

\newpage
\textbf{3. Genuine variables, identical laws and fourth moments.}
Put $U_i=\lambda(g_i)$ and $X_i=U_i+U_i^*$. These are the images of the
actual factor-group-algebra element $g+g^{-1}$. Their self-adjointness
and membership in $\A_i$ are proved, and $U_i^3=1$ gives
\[
 X_i^2=X_i+2,\qquad X_i^3=3X_i+2,\qquad X_i^4=5X_i+6.
\]
Since neither $g$ nor $g^{-1}$ is the identity of $C_3$, the true
coefficient state gives $\tau(X_i)=0$. Applying the state to these
identities yields the actual values
\[
 \tau(X_i^2)=2,\qquad \tau(X_i^3)=2,\qquad \tau(X_i^4)=6.
\]
No moment is assigned as a definition. The identities are obtained in
the native factor group algebra, mapped through the genuine bounded
representation, and evaluated by the actual state.

Define $Y_i=X_i/\sqrt2$. Real scalar multiplication preserves the
proved self-adjointness and factor membership. Direct normalization
gives $\tau(Y_i)=0$, $\tau(Y_i^2)=1$, and $\tau(Y_i^4)=3/2$, with the
explicit uniform bound $\|\tau(Y_i^4)\|\le2$. The full polynomial law
is also checked: for every $p\in\C[z]$ and all $i,j$,
\[
 \tau\bigl(p(Y_i)\bigr)=\tau\bigl(p(Y_j)\bigr).
\]
Indeed evaluation commutes with the actual factor-algebra homomorphism,
and the trace of any factor image is its original identity coefficient.
Since each variable is self-adjoint, this proves equality of their
one-variable star-polynomial distributions. Together with the full
factor-algebra freeness above, the sequence is free i.i.d.

\textbf{4. The actual normalized sum and its third moment.}
The Lean development forms the genuine finite sums of these bounded
operators, and also the same sums in the native star-subalgebra $\A$.
It proves that their operator images agree and that
\[
 S_n=\frac1{\sqrt n}\sum_{i<n}Y_i
     =\frac1{\sqrt{2n}}\sum_{i<n}X_i.
\]
All square roots are native real square roots, embedded in $\C$.
For actual indices $i,j,k$, the exact mixed third moment is
\[
 \tau(X_iX_jX_k)=
 \begin{cases}2,&i=j=k,\\0,&\text{otherwise.}\end{cases}
\]
If adjacent indices differ, full freeness applies. If exactly two
adjacent indices agree, replace their product by $X_i+2$ and use
centering and the two-factor freeness identity. The all-equal case
uses the genuine cubic identity. Expanding the true cube of the finite
sum and evaluating all three finite sums therefore gives
\[
 \tau\left(\left(\sum_{i<n}X_i\right)^3\right)=2n.
\]
For $n>0$, the proved square-root identity $(\sqrt{2n})^2=2n$
and its nonzero value imply
\[
 \tau(S_n^3)=\frac{2n}{(\sqrt{2n})^3}
            =\frac1{\sqrt{2n}}.
\]
The final rate statement concerns this actual state of the actual
operator sum; a separate sequence with a stipulated value is never
substituted.

\textbf{5. The actual semicircle target.}
Lean defines the real measure
\[
 d\mu_{\rm sc}(x)=\mathbf1_{[-2,2]}(x)
       \frac{\sqrt{4-x^2}}{2\pi}\,dx
\]
by native restriction of Lebesgue measure and \texttt{withDensity}.
This is the standard variance-one semicircle density
\cite[p.~327]{Novak}, with the normalization matching the $Y_i$ above.
Its normalization is proved, rather than assumed: the change of
variables $x=2t$ reduces its integral to the native theorem
\[
 \int_{-1}^{1}\sqrt{1-t^2}\,dt=\frac\pi2.
\]
The resulting instance is a true \texttt{IsProbabilityMeasure}.
The density is continuous and nonnegative. Continuity on the actual
compact support also proves the cubic integrand is genuinely integrable
against this measure. Finally, its weighted cubic integrand is odd;
native interval reflection and linearity give
$\int x^3\,d\mu_{\rm sc}=0$. This is not a default value of an
undefined or nonintegrable Bochner integral.

\newpage
\textbf{6. Quantified failure of the claimed rate.}
Let
\[
 E_n=\tau(S_n^3)-\int x^3\,d\mu_{\rm sc}(x)\in\C.
\]
The two actual calculations above prove $\|E_n\|=1/\sqrt{2n}$
for every $n>0$. Given any $C\ge0$ and starting index $N$, choose a
natural number $n\ge N$ with $n>0$ and $n>2C^2$. Then
\[
 (C\sqrt{2n})^2=2C^2n<n^2,
 \qquad
 \frac Cn<\frac1{\sqrt{2n}}=\|E_n\|.
\]
All factors in the squaring and division steps are proved nonnegative
or positive. Lean proves the full quantifier statement
\[
 \forall C\ge0\ \forall N\in\N\ \exists n\ge N:
 \quad n>0\quad\hbox{and}\quad C/n<\|E_n\|.
\]
The native eventual-bound characterization then yields
\[
 \neg\,\operatorname{Asymptotics.IsBigO}_{n\to\infty}
       \bigl(E_n,n^{-1}\bigr).
\]
Thus the first universal rate clause in the printed conjecture fails
already for scalar-valued, centered, bounded, self-adjoint, free i.i.d.
variables with unit variance and uniformly bounded fourth moments.
The usual free central limit theorem remains compatible with this
example: the third-moment discrepancy does tend to zero. The
$O(n^{-1})$ bound asserted without an odd-moment hypothesis is the
specific obstruction established here.

\textbf{Formal objects and reproduction.}
The complete final theorem is \texttt{TLMC7730.complete\_counterexample}.
The native countable group coproduct, unique reduced words,
$\ell^2$ Hilbert space, actual bounded regular representation, native
image star-subalgebra, bounded positive unital tracial state, whole
factor subalgebras and all-length freeness, polynomial laws, real finite
sums, actual semicircle probability measure and native asymptotic
predicate are connected in the same development. The general free CLT
is not needed as a premise for the finite-sum moment calculations or
the rate contradiction.

The project pins Lean~4.33.0 and Mathlib commit
\begin{center}
{\footnotesize\texttt{db584cd6d46c92f209a44c0f1c829460d327499d}}.
\end{center}
From \texttt{lean/}, run the following four commands separately:
\begin{verbatim}
lean --version
lake build +Main
lake env lean -DwarningAsError=true Check.lean
lake env leanchecker --fresh Main
\end{verbatim}
\texttt{lean/verify.sh} runs the last three; the version check is separate.
The audit calls ordinary \texttt{\#print axioms} on every owned declaration
in the actual imported environment, including generated names. Only
\texttt{propext}, \texttt{Classical.choice} and \texttt{Quot.sound} are
permitted. The submission retains raw source-bound command logs,
execution receipts, actual owned-thread masks and a SHA-256 manifest.
Production durations are diagnostic and supply no controlled
efficiency conclusion. Organizer review remains pending.

\begin{thebibliography}{1}\small
\bibitem{Novak} J.~Novak, with illustrations by M.~LaCroix,
\emph{Three lectures on free probability}, in \emph{Random Matrices},
MSRI Publications~65 (2014), 309--383.
Relevant definitions and models: pp.~327, 329--330, 332, 339--340.
\href{https://library.slmath.org/books/Book65/files/140819-Novak.pdf}{Publisher-hosted full text}.
\end{thebibliography}
\end{document}
